% Generated draft: requires mathematical and source review.
\documentclass[11pt]{article}
\usepackage[margin=25mm]{geometry}
\usepackage{fontspec}
\setmainfont{Latin Modern Roman}
\usepackage{amsmath,amssymb,mathtools,unicode-math}
\setmathfont{Latin Modern Math}
\usepackage{enumitem,needspace,xurl,hyperref}
\setlength{\parindent}{0pt}
\setlength{\parskip}{5pt}
\setlength{\emergencystretch}{3em}
\newcommand{\N}{\mathbb{N}}
\newcommand{\Z}{\mathbb{Z}}
\newcommand{\Q}{\mathbb{Q}}
\newcommand{\R}{\mathbb{R}}
\newcommand{\C}{\mathbb{C}}
\newcommand{\F}{\mathbb{F}}
\newcommand{\A}{\mathbb{A}}
\newcommand{\PP}{\mathbb{P}}
\newcommand{\E}{\mathbb{E}}
\newcommand{\eps}{\varepsilon}
\newcommand{\dd}{\,\mathrm{d}}
\begin{document}
\section*{UnsolvedMath ID: 30003110}
\textbf{Source problem code:} catalogue entry
\subsection{Definitions and mathematical statement}
Let \(K\) be a fixed number field. For an elliptic curve \(E/K\), write
\[
r_K(E)=\operatorname{rank}_{\Z}\bigl(E(K)/E(K)_{\mathrm{tors}}\bigr).
\]
Choose a basis \(P_1,\ldots,P_r\) of this free abelian group, where \(r=r_K(E)>0\). Using the N\'eron--Tate height \(\widehat h_E\) associated with the origin divisor, fix the pairing convention
\[
\langle P,Q\rangle_E=\frac{\widehat h_E(P+Q)-\widehat h_E(P)-\widehat h_E(Q)}{2},
\qquad
\operatorname{Reg}_K(E)=\det\bigl(\langle P_i,P_j\rangle_E\bigr)_{1\leq i,j\leq r}.
\]
The determinant is independent of the chosen integral basis. The supplied record does not specify its height normalization; the displayed convention makes the regulator used here explicit.

For every integer \(R\geq1\) and real number \(B>0\), put
\[
\mathcal E_K(R,B)=\{E/K:\ E\text{ is an elliptic curve},\ 1\leq r_K(E)\leq R,\ \operatorname{Reg}_K(E)\leq B\}.
\]
The intended conditional finiteness assertion has the following form, once the hypothesis and equivalence relation below are clarified:
\[
\text{Assuming ABC},\qquad
\forall R\in\Z_{\geq1}\ \forall B\in\R_{>0},\qquad
\#\bigl(\mathcal E_K(R,B)/{\sim}\bigr)<\infty.
\]
Here the source explicitly calls \(\sim\) ``\(\Q\)-isomorphism.'' For arbitrary elliptic curves defined over \(K\), this does not specify whether one means isomorphism over \(K\), isomorphism after forgetting to \(\Q\)-schemes, or isomorphism of models defined over \(\Q\). The last interpretation would require a descent hypothesis absent from the record. Consequently the displayed assertion is a formulation template, not a fully specified conjecture; no equivalence relation has been silently substituted. The source also does not specify whether ABC means its rational or number-field formulation, or precisely which fields its hypothesis must cover.

The older \texttt{original\_statement} additionally records a Lang--Silverman conditional finiteness assertion for simple abelian varieties over a fixed number field, with fixed dimension, bounded positive Mordell--Weil rank, and bounded regulator, again using the phrase \(\Q\)-isomorphism classes. The supplied verification says that this duplicated assertion belongs to \texttt{OWR-14603-007} and was removed from the present target. The present record therefore targets only the ABC consequence for elliptic curves. Its supplied literature assessment still discusses the Lang--Silverman hypothesis; that assessment does not resolve the formulation issues above.
\subsection{Short English statement}
Assuming the intended version of ABC, the current target asks whether elliptic curves over a fixed number field, with positive rank bounded above and regulator bounded above, have only finitely many classes. The source's meaning of \(\Q\)-isomorphism and its version of ABC require clarification. The supplied open label and conditional-result descriptions are source assertions, not independent mathematical verification.
\subsection{Sources}
\begin{itemize}
\item Local registry: \texttt{data/unsolvedmath/problems.json}, supplied record with permanent UnsolvedMath ID \(30003110\), exact source code \texttt{catalogue entry}, title ``Regulator Finiteness for Elliptic Curves from ABC.'' No individual proposer is supplied; the proposed year is recorded as 2016.
\item Supplied report citation: \emph{Diophantische Approximationen (2017)}, Oberwolfach Reports, DOI \texttt{10.4171/owr/2016/21}, pp.\ 1099--1187. \url{https://doi.org/10.4171/owr/2016/21}. The report title/year and DOI are reproduced as supplied.
\item Supplied literature source, ``OWR report'': \url{https://ems.press/content/serial-article-files/46625?nt=1}.
\item Supplied literature source, ``Silverman notes'': \url{https://swc-math.github.io/aws/2024/2024SilvermanNotes.pdf}.
\item Literature assessment dated exactly \texttt{2026-08-21}: supplied dated curation, not independently verified. It labels the problem open and describes the recorded theorem as conditional on Lang--Silverman, with special consequences established but the hypothesis unresolved. The supplied background explicitly requests independent verification of time-sensitive status and cited claims.
\item The supplied statement status is \texttt{corrected\_verified}. Its verification note reports that the official report was checked and that the duplicated Lang--Silverman consequence was removed in favor of the distinct ABC consequence for elliptic curves. This reports the registry's verification claim; no source visit or independent verification is asserted here.
\end{itemize}
\end{document}
